% ATTEMPT_STATUS: unresolved
% TOP_PROBLEM: {"problemId": "problem.p-versus-np", "canonicalTitle": "P versus NP", "releaseRank": 1, "primaryDomain": "theoretical_computer_science", "primaryDomainLabel": "Theoretical computer science", "displayStatus": "open", "releaseStatus": "open"}
% !TeX program = lualatex
% Complete problem section copied verbatim from research_batch01.tex.
\documentclass[11pt]{article}
\usepackage[margin=25mm]{geometry}
\usepackage{fontspec}
\setmainfont{Latin Modern Roman}
\usepackage{amsmath,amssymb,mathtools}
\usepackage{unicode-math}
\setmathfont{Latin Modern Math}
\usepackage{enumitem,needspace}
\usepackage{xurl,hyperref}
\hypersetup{colorlinks=true,urlcolor=blue}
\setcounter{secnumdepth}{1}
\setlength{\parindent}{0pt}
\setlength{\parskip}{5pt}
\setlength{\emergencystretch}{3em}
\setlist[itemize]{leftmargin=3em,itemsep=5pt,topsep=4pt,parsep=0pt}
\allowdisplaybreaks
\newcommand{\N}{\mathbb{N}}
\newcommand{\Z}{\mathbb{Z}}
\newcommand{\Q}{\mathbb{Q}}
\newcommand{\R}{\mathbb{R}}
\newcommand{\C}{\mathbb{C}}
\newcommand{\F}{\mathbb{F}}
\newcommand{\A}{\mathbb{A}}
\newcommand{\PP}{\mathbb P}
\newcommand{\E}{\mathbb{E}}
\newcommand{\eps}{\varepsilon}
\newcommand{\dd}{\,\mathrm d}
\newcommand{\sref}[2]{\hyperlink{source:#1:#2}{[S#2]}}
\newcommand{\eref}[2]{\hyperlink{extra:#1:#2}{[E#2]}}
% Turn the next switch off for a shorter reading copy. The quotations preserve
% every exactTarget field, including qualifications omitted from a short summary.
\newif\ifshowcatalogue
\showcataloguetrue
\newcommand{\cataloguescope}[1]{\ifshowcatalogue
 \paragraph{Exact scope in the supplied catalogue.}
 \begin{quote}\small #1\end{quote}\fi}
\begin{document}
\setcounter{section}{0}
% ================================================================
% problemId: problem.p-versus-np
% source releaseRank: 1; source releaseStatus: open
% exactTarget SHA-256: 9ac81ed7e9ada6c6d9d77326bebb48362f3a3a5574e196911b79e5cc09752960
\clearpage
\section{\texorpdfstring{P versus NP}{P versus NP}}\label{1}
\subsection{Definitions and mathematical statement}
A language is a set $L\subseteq\{0,1\}^*$. Write $\mathsf P$ for languages decided by a deterministic Turing machine in time bounded by a polynomial in the input length. A language belongs to $\mathsf{NP}$ if there are a deterministic polynomial-time verifier $V$ and a polynomial $p$ such that
\[
 x\in L\quad\Longleftrightarrow\quad \exists y\in\{0,1\}^{\le p(|x|)}\;V(x,y)=1.
\]
The polynomial and machine may depend on $L$, but not on its individual inputs. The question is whether $\mathsf P=\mathsf{NP}$. A negative answer requires a language in $\mathsf{NP}$ outside $\mathsf P$, not merely a slow algorithm for one problem.
\par\smallskip\textit{Formulation and concept references:} \sref{1}{1}, \eref{1}{1}.
\cataloguescope{Determine whether P = NP, that is, whether every decision problem whose solutions can be verified in polynomial time by a deterministic Turing machine can also be solved in polynomial time by a deterministic Turing machine.}
\subsection{Short English statement}
Does an efficient way to check a proposed yes-answer always imply an efficient way to decide whether the answer is yes?
% BEGIN RESEARCH ATTEMPT -- batch 1, problem 1, 24 September 2026
\subsection{Research attempt}

\paragraph{Current frontier and source audit.}
The target here is the ordinary, uniform, classical language-class question.
Cook's NP-completeness theorem makes a deterministic polynomial-time algorithm
for Boolean circuit satisfiability sufficient for equality; a lower bound for
one restricted representation of satisfying assignments is not sufficient for
inequality \eref{1}{1}. An established milestone on the lower-bound side is
Williams's nonuniform ACC lower bound for nondeterministic exponential time,
not a lower bound for arbitrary circuits computing an NP language
\eref{1}{5}. As a recent, model-specific result, the March 2026 manuscript of
Chen, Tal and Wang reports $n^{5/2-\varepsilon}$ gate lower bounds against
depth-two threshold circuits for a function in $\mathsf E^{\mathsf{NP}}$
\eref{1}{6}. The class, gate basis and quantifiers in that statement cannot be
replaced by those of this problem. These are relevant checkpoints, not a claim
that one numerical bound describes the entire frontier.

Three obstructions must be kept separate. Baker--Gill--Solovay construct oracle
worlds with opposite answers; hence an argument valid unchanged for every
oracle cannot settle this question \eref{1}{2}. The natural-proofs obstruction
is conditional on sufficiently strong pseudorandomness and concerns its
specific constructivity, largeness and usefulness requirements; it is not an
unconditional prohibition on combinatorial lower bounds \eref{1}{3}.
Aaronson--Wigderson's algebraic-oracle results show why merely replacing Boolean
queries by low-degree extensions does not supply a general escape
\eref{1}{4}. None of these results says that an explicit, syntax-sensitive
argument is impossible.

The existing catalogue claims were also checked for scope. The abstract of
Hole's September 2026 version explicitly uses a finiteness condition and a
constructive interpretation \sref{1}{3}. Under the present definition, however,
$\mathsf{NP}\subseteq\mathsf{EXPTIME}$ follows directly: on input length $n$,
try all fewer than $2^{p(n)+1}$ permitted witnesses and run the fixed verifier
on each. This argument does not require a theory that proves an algorithm
correct. Thus an assertion about unprovability of algorithmic correctness is
not by itself a separation of these language classes. The PNP Labs page
contains self-reported formalized components together with explicitly
unproved global and polynomial-runtime endpoints \sref{1}{4}; those components
are not used as a theorem settling the target.

\paragraph{Attack map and selected route.}
A diagonalization route would have to diagonalize against every polynomial
exponent while retaining one polynomially bounded NP verifier. Simulating the
$i$th machine for $n^i$ steps does not accomplish that: the exponent then varies
with the input. A circuit route would need a superpolynomial lower bound for
an NP language against unrestricted circuits; obtaining that would additionally
settle Section~\ref{12} negatively, a stronger conclusion than is requested
here. An upper-bound route can instead exploit the explicit gate syntax of a
verifier and try to eliminate existential variables while retaining a compact
exact description of the remaining feasible set. The last route is pursued
below. Its precise missing ingredient is a \emph{uniformly computable,
polynomial-size summary preserved through all the required eliminations},
not merely a succinct input formula. The distinction between succinctness,
tractable queries and tractable transformations is central in knowledge
compilation \eref{1}{7}.

The proposed combination is linear algebra over $\F_2$, affine summaries of
nonlinear gates, and decomposition into independent factors. The reasoning
below was developed for this notebook. The elementary lemmas are proved here;
no claim of historical novelty or of a new general complexity lower bound is
made.

\paragraph{Attempt: an exact algebraic encoding.}
For a circuit with AND, OR and NOT gates, introduce a variable for every input
and every gate output. Over $\F_2$, impose, respectively,
\begin{equation}\label{ra:1:gates}
 z+xy=0,\qquad z+x+y+xy=0,\qquad z+1+x=0,
\end{equation}
and impose $z_{\rm out}=1$. All variables range over $\F_2$; no relaxation to
arbitrary field extensions is being made. Induction in a topological ordering
of the circuit proves that every input has exactly one extension to its gate
values, and that this extension satisfies the output equation exactly when
the circuit accepts. There are polynomially many variables and quadratic
equations in the circuit's bit length. Conversely, a proposed solution of such
a system is polynomially checkable. Consequently a uniform polynomial-time
solver for all these quadratic systems would decide circuit satisfiability
and hence give $\mathsf P=\mathsf{NP}$ \eref{1}{1}.

Gaussian elimination solves affine systems and computes coordinate projections
of affine subspaces. This suggests storing a feasible relation as a union of
affine subspaces, and using linear algebra on each piece. The following sharp
calculation tests how much such a representation can compress even a quadratic
graph.

\medskip
\textbf{Lemma 1.1 (exact affine-cover size of a quadratic graph).}
Let $Q:\F_2^d\to\F_2^m$ have coordinate polynomials of degree at most two, in
multilinear form, and define its vector-valued polar form by
\[
 B(u,v)=Q(u+v)+Q(u)+Q(v)+Q(0).
\]
This is alternating and bilinear. Put
\[
 \alpha(B)=\max\{\dim U:U\le\F_2^d,\ B(u,v)=0\text{ for all }u,v\in U\},
 \qquad \Gamma_Q=\{(x,Q(x)):x\in\F_2^d\}.
\]
If $\kappa(\Gamma_Q)$ is the least number of affine subspaces whose union is
\emph{exactly} $\Gamma_Q$, then
\begin{equation}\label{ra:1:cover}
             \kappa(\Gamma_Q)=2^{d-\alpha(B)}.
\end{equation}
The same minimum holds when each piece may instead be the coordinate
projection of an affine subspace with any number of auxiliary variables.

\emph{Proof.}
Let $A\subseteq\Gamma_Q$ be an affine subspace. Projection onto the input
coordinates is injective on $A$, since a graph has only one output per input.
Its image is an affine space $x_0+U$ of dimension $\dim A$. The restriction of
$Q$ to that space is affine. Its second additive difference is therefore zero:
\[
 Q(x_0)+Q(x_0+u)+Q(x_0+v)+Q(x_0+u+v)=B(u,v)=0
 \quad(u,v\in U).
\]
Thus $\dim A\le\alpha(B)$ and $|A|\le2^{\alpha(B)}$. Since the graph has $2^d$
points, any exact cover has at least $2^{d-\alpha(B)}$ pieces, even with overlap.

For the reverse bound, choose a totally isotropic $U$ of maximum dimension.
On every coset $x_0+U$, the map
\[
 u\longmapsto Q(x_0+u)+Q(x_0)
\]
is additive by the displayed second-difference identity, and hence linear over
$\F_2$. The graph above that coset is an affine space. The $2^{d-\alpha(B)}$
cosets partition the input space and their graphs partition $\Gamma_Q$.
Finally a coordinate projection of an affine subspace is itself affine.
An exact union cannot contain a piece with a point outside the graph, so
auxiliary affine variables cannot evade the same lower bound.\hfill$\square$

\medskip
\textbf{Corollary 1.2 (a constructive special-case solver).}
Given $Q$, $b\in\F_2^m$, and an explicitly supplied basis of a totally isotropic
subspace $U$ of dimension $r$, solvability of $Q(x)=b$ can be decided in
\begin{equation}\label{ra:1:algorithm}
             2^{d-r}\operatorname{poly}(d,m)
\end{equation}
bit operations, with polynomial working space when solutions need not all be
stored. Indeed extend the basis of $U$ to a basis of $\F_2^d$, enumerate the
$2^{d-r}$ complementary representatives $x_0$, and solve the linear system
\[
 Q\!\left(x_0+\sum_{j=1}^r t_ju_j\right)
   =Q(x_0)+\sum_{j=1}^r t_j\bigl(Q(x_0+u_j)+Q(x_0)\bigr)=b.
\]
All coefficients are bits. A polynomial bound for the displayed enumeration
would follow from $d-r=O(\log N)$, where $N$ is the encoded input size,
\emph{provided the required basis is also found in polynomial time}. Neither
the existence nor the efficient discovery of such a basis is assumed for a
general circuit encoding. When the constraint system is a conjunction of
independent systems on disjoint variable sets, the work is a \emph{sum} of
these costs over the components, not their product. This is the useful part
of combining affine algebra with factorization.

\paragraph{Attempt: falsifying the proposed compression invariant.}
Take
\[
 Q_n(x_1,y_1,\ldots,x_n,y_n)=(x_1y_1,\ldots,x_ny_n).
\]
For the $i$th coordinate of its polar form,
\[
 B_i(u,v)=u_{x_i}v_{y_i}+u_{y_i}v_{x_i}.
\]
If $U$ is totally isotropic, its projection onto each pair $(x_i,y_i)$ has
dimension at most one: a two-dimensional projection contains two vectors
with determinant one, contradicting isotropy. The combined coordinate
projections are injective, so $\dim U\le n$. Equality is attained by $y_i=0$
for all $i$. Lemma 1.1 gives
\begin{equation}\label{ra:1:andcover}
                        \kappa(\Gamma_{Q_n})=2^n.
\end{equation}
Nevertheless its graph is described by just $3n$ CNF clauses,
\[
 \bigwedge_{i=1}^n
  \bigl[(\neg z_i\vee x_i)\wedge(\neg z_i\vee y_i)
        \wedge(z_i\vee\neg x_i\vee\neg y_i)\bigr],
\]
with $O(n\log(n+1))$ bits in a standard variable-index encoding. The all-zero
assignment satisfies it, and membership in the graph is checked in linear
many Boolean operations. Even satisfiability after arbitrary variable
restrictions is decided independently on the $n$ triples. Thus a universal
polynomial bound on flat affine summaries is \emph{false on a plainly
tractable family}. The exact cover bound cannot be promoted to a lower bound
for arbitrary SAT algorithms.

Factorization repairs this particular example: keep the $n$ ternary relations
as independent factors rather than distributing their conjunction over all
affine choices. But connectedness alone is not the missing invariant. The
single scalar equation
\[
                  q(x,y)=\sum_{i=1}^n x_i y_i=b
\]
involves every variable and has a nondegenerate symplectic polar form. If $U$
is isotropic, $U\subseteq U^\perp$ and $\dim U^\perp=2n-\dim U$, giving
$\alpha(B)=n$. Its graph again needs $2^n$ affine pieces. Yet $b=0$ is solved
by the zero vector and $b=1$ by $x_1=y_1=1$, all other coordinates zero.
Large affine-cover complexity, even for a connected quadratic relation, need
not imply decision hardness.

\paragraph{Attempt: can quadratic factors be retained through projection?}
A second possible repair is to store a conjunction of quadratic equations,
without flattening it, and eliminate internal variables while remaining in
that language. There is already a finite obstruction. Eliminating $w$ from
\[
                     w=xy,\qquad z=wt
\]
gives the four-variable relation
\begin{equation}\label{ra:1:cubic}
                   R=\{(x,y,t,z):z=xyt\}.
\end{equation}
It cannot be defined by \emph{any} conjunction of quadratic equations over
$\F_2$ on these four visible variables alone.

To prove this, reduce every polynomial modulo $v^2-v$ in each variable; this
does not increase degree or change its Boolean zero set. A quadratic
polynomial vanishing on $R$ has the form
\[
 q_0(x,y,t)+z\ell(x,y,t),\qquad \deg q_0\le2,\quad\deg\ell\le1.
\]
For the seven triples different from $(1,1,1)$, the graph has $z=0$, so $q_0$
vanishes there. A Boolean function supported only at $(1,1,1)$ is either zero
or $xyt$: the squarefree monomials form a basis of all functions on
$\F_2^3$. As $q_0$ has degree at most two, it must be the zero function.
Evaluating at the remaining graph point then gives $\ell(1,1,1)=0$.
Consequently \emph{every} quadratic vanishing on $R$ also vanishes at
$(1,1,1,0)$, which is not in $R$. No number of such equations can cut out $R$.
This disproves the proposed closure lemma, rather than merely failing to prove
an efficient implementation.

Allowing degree three represents this example immediately; keeping $w$ as an
existential variable also represents it. The former permits degree growth
under repeated elimination, and the latter has not eliminated the search
problem. Neither observation rules out more flexible algorithms, but neither
closes the proposed upper bound.

\paragraph{Computational check.}
The accompanying script \texttt{research\_batch01\_checks.py} enumerates every
linear subspace of $\F_2^{2n}$ in reduced row-echelon form for $n=1,2,3$,
checks isotropy on basis pairs, and independently forms the affine graph
pieces over their cosets. The total subspace counts are $5,67,2825$; the
maximum isotropic dimensions are $1,2,3$, with $3,9,27$ maximum-dimensional
isotropic subspaces. The corresponding largest graph-piece sizes are $2,4,8$
in graphs of sizes $4,16,64$. Exact counting and the explicit partition give
cover numbers $2,4,8$ in those cases. The script also examines all $2^{11}$
multilinear degree-at-most-two polynomials on four variables. Exactly eight
vanish on \eqref{ra:1:cubic}; their common zero set has nine points, including
the spurious point $(1,1,1,0)$. These are finite exact checks of the local
arguments, not experimental evidence for either answer to P versus NP.

\paragraph{Stress test and dependencies.}
The lower bound counts representations of an \emph{entire relation}, whereas
the target permits an algorithm that only decides nonemptiness. The easy
families above invalidate the inference from the former task to the latter.
The isotropic-space maximum in Lemma 1.1 is existential; its proof is not a
polynomial-time algorithm for finding that maximum. The quadratic-projection
obstruction concerns a specified output language, not all data structures or
Turing machines. These distinctions are instances of the compilation issues
in \eref{1}{7}, and arise before applying the general barriers in
\eref{1}{2}, \eref{1}{3}, \eref{1}{4}.

A stronger, still unproved route would have to give a syntax-sensitive
elimination or decomposition algorithm whose \emph{total} intermediate bit
length and total arithmetic work are polynomial for every gate system
\eqref{ra:1:gates}. It would need to handle coupled nonlinear factors without
flat affine expansion, without unbounded-degree expansion, and without
hiding a satisfiability oracle inside a summary operation. The two simplest
universal closure claims tested here are false; no substitute with these
properties has been established.

Were that algorithm supplied, it would imply $\mathsf P=\mathsf{NP}$ and
therefore $\mathsf{NP}=\mathsf{coNP}$ in Section~\ref{13}. It would also rule
out the one-way functions of Section~\ref{8}: for a given output $y$ and
length $n$, the existence of a preimage extending any chosen prefix is an NP
predicate; polynomially many prefix-extension decisions recover a preimage
whenever one exists. This implication uses the exact length and any-preimage
conventions of that entry. Conversely, failure of the two proposed
representations establishes neither $\mathsf P\ne\mathsf{NP}$ nor existence
of cryptographic hardness.

\paragraph{Outcome.}
\textbf{Outcome: Exploratory approach with unresolved gap.}

The obtained mathematics is the exact quadratic-graph affine-cover formula,
a conditional-on-a-supplied-subspace solver, and two explicit obstructions to
natural elimination schemes. These are proved within the restricted models
stated above; their historical novelty has not been established. The attempt
has not found either a general polynomial-time SAT algorithm or a
superpolynomial lower bound for all such algorithms. Its useful negative
conclusion is specific: small verifier syntax does not force small flat
affine summaries, and conjunctions of visible quadratic equations are not
closed under existential projection. A successful continuation must change
the representation or the invariant, not merely sharpen constants in either
failed claim.
% END RESEARCH ATTEMPT -- batch 1, problem 1

\subsection{Sources}
\begin{itemize}\raggedright
\item[\textnormal{[S1]}]\hypertarget{source:1:1}{} Clay Mathematics Institute, The Millennium Prize Problems.
\url{https://www.claymath.org/library/monographs/MPPc.pdf}.
{\small\textit{Catalogue formulation source.}}
\item[\textnormal{[S2]}]\hypertarget{source:1:2}{} P vs NP.
\url{https://www.claymath.org/millennium/p-vs-np/}.
{\small\textit{Catalogue research/status reference; not a proof endorsement.}}
\item[\textnormal{[S3]}]\hypertarget{source:1:3}{} Constructive solvability and the P versus NP problem — Arne Hole.
\url{https://arxiv.org/abs/2406.16843}.
{\small\textit{Catalogue research/status reference; not a proof endorsement.}}
\item[\textnormal{[S4]}]\hypertarget{source:1:4}{} PNP Labs — P versus NP formal reconstruction.
\url{https://pnplabs.com.au/}.
{\small\textit{Catalogue research/status reference; not a proof endorsement.}}
\item[\textnormal{[E1]}]\hypertarget{extra:1:1}{} Stephen Cook, The P versus NP Problem, Clay Mathematics Institute, Section 1 and Appendix.
\url{https://www.claymath.org/wp-content/uploads/2022/06/pvsnp.pdf}.
{\small\textit{Added primary source: Definitions and formulation. Consulted 24 September 2026.}}

% BEGIN ADDED RESEARCH SOURCES -- batch 1, problem 1
\item[\textnormal{[E2]}]\hypertarget{extra:1:2}{} Theodore Baker, John Gill and
Robert Solovay, \emph{Relativizations of the $\mathcal P=?\mathcal{NP}$ Question},
SIAM Journal on Computing 4(4) (1975), 431--442, oracle equality and separation
results. \url{https://doi.org/10.1137/0204037}.
{\small\textit{Consulted 24 September 2026: publisher abstract; used only for
the stated oracle obstruction, not a new proof audit.}}
\item[\textnormal{[E3]}]\hypertarget{extra:1:3}{} Alexander A. Razborov and Steven
Rudich, \emph{Natural Proofs}, Journal of Computer and System Sciences 55(1)
(1997), 24--35; author manuscript dated 30 November 1996, abstract and the
natural-property/hardness framework.
\url{https://www.cs.umd.edu/~gasarch/BLOGPAPERS/natural.pdf}.
{\small\textit{Consulted 24 September 2026. The pseudorandomness assumption is
retained; no unconditional barrier is asserted.}}
\item[\textnormal{[E4]}]\hypertarget{extra:1:4}{} Scott Aaronson and Avi Wigderson,
\emph{Algebrization: A New Barrier in Complexity Theory}, ACM Transactions on
Computation Theory 1(1) (2009), Article 2; author manuscript, Sections 1--2
and 5, especially Theorems 5.1 and 5.3.
\url{https://www.scottaaronson.com/papers/alg.pdf}.
{\small\textit{Consulted 24 September 2026: definitions and barrier statements.}}
\item[\textnormal{[E5]}]\hypertarget{extra:1:5}{} Ryan Williams,
\emph{Nonuniform ACC Circuit Lower Bounds}, Journal of the ACM 61(1) (2014),
Article 2, 32 pages; author manuscript, Theorem 1.1 and Section 1.1.
\url{https://people.csail.mit.edu/rrw/acc-lbs.pdf}.
\url{https://doi.org/10.1145/2559903}.
{\small\textit{Consulted 24 September 2026.}}
\item[\textnormal{[E6]}]\hypertarget{extra:1:6}{} Lijie Chen, Avishay Tal and
Yichuan Wang, \emph{Super-quadratic Lower Bounds for Depth-2 Linear Threshold
Circuits}, ECCC TR26-039, 15 March 2026, abstract and reported
$\mathsf E^{\mathsf{NP}}$ threshold-circuit bound.
\url{https://eccc.weizmann.ac.il/report/2026/039/}.
{\small\textit{Consulted 24 September 2026: report page and abstract. A
recent manuscript result, not an independent verification of its full proof.}}
\item[\textnormal{[E7]}]\hypertarget{extra:1:7}{} Adnan Darwiche and Pierre
Marquis, \emph{A Knowledge Compilation Map}, Journal of Artificial Intelligence
Research 17 (2002), 229--264, Sections 1--2 and the separation of succinctness,
queries and transformations.
\url{https://arxiv.org/abs/1106.1819}.
\url{https://doi.org/10.1613/jair.989}.
{\small\textit{Consulted 24 September 2026: article introduction and language
definitions. Used as methodological context, not as the source of Lemma 1.1.}}
% END ADDED RESEARCH SOURCES -- batch 1, problem 1
\end{itemize}

\end{document}
